
\section*{\underline{EXPOSÉ IX}\\[3pt] \underline{\underline{QUELQUES CALCULS DE GROUPES $K$}}}
\begin{center}
\underline{par P. BERTHELOT}
\end{center}

Cet exposé ne correspond à aucun exposé oral, et ne sera pas utilisé dans la suite du séminaire. Il contient un certain nombre de calculs de groupes $K_\bullet$ pour des fibrations diverses : fibrés vectoriels, fibrés projectifs, torseurs sous un groupe linéaire, etc\ldots{} Dans tout l'exposé, le schéma de base sera supposé noethérien, et la notation $K_\bullet(X)$ désignera le groupe de Grothendieck des faisceaux cohérents sur un schéma $X$. On rappelle que si $X$ est noethérien, l'homomorphisme du $K$ de la catégorie des $\underline{O}_X$-Modules cohérents dans le $K$ de la catégorie des complexes pseudo-cohérents sur $X$ est un isomorphisme (IV 2.4).

Les résultats exposés ici sont l'analogue de ceux que l'on obtient pour l'anneau de Chow, et qu'on trouvera dans les exposés de Grothendieck au Séminaire Chevalley 1958. La question se pose également de déterminer dans quelle mesure des énoncés du même type peuvent être obtenus pour le gradué associé au $K_\bullet$ pour la filtration définie par la dimension du support (cf. X). Il ne semble pas que la réponse en soit connue, même pour la proposition 1.1 (suite exacte d'homotopie) qui est l'une des propriétés les plus utilisées dans les calculs qui suivent.

\subsection*{1. Fibrés vectoriels}

\textbf{Proposition 1.1. (Suite exacte d'homotopie).} \uline{Soient} $X$ \uline{un schéma noethérien}, $X'$ \uline{un sous-schéma fermé de} $X$, $U$ \uline{l'ouvert complémentaire de} $X'$ \uline{dans} $X$, $i$ \uline{l'immersion de} $X'$ \uline{dans} $X$, \uline{et} $j$ \uline{l'immersion de} $U$ \uline{dans} $X$. \uline{Alors la suite} :
